%% =====================================================================
%%  Pages liminaires communes — cahiers OptiTAB Maternelle
%% =====================================================================
% Page de titre
\thispagestyle{empty}
\begin{tikzpicture}[remember picture,overlay]
\begin{scope}[shift={(current page.south west)}]
  \node[anchor=north,inner sep=0pt] at (105mm,262mm) {\includegraphics[width=62mm]{logo-couleur}};
  \node[anchor=base,inner sep=0pt,text=matAccent,font=\otSemi\fontsize{10}{12}\selectfont] at (105mm,186mm) {\otCaps[22]{Maternelle · \matNiveau}};
  \node[anchor=base,inner sep=0pt,text=otNavy,font=\otBold\fontsize{40}{44}\selectfont] at (105mm,166mm) {Mathématiques};
  \node[anchor=base,inner sep=0pt,text=otInk,font=\otLight\fontsize{22}{26}\selectfont] at (105mm,151mm) {Mon cahier d’activités};
  \fill[matAccent] (93mm,141mm) rectangle (117mm,142.2mm);
  \node[anchor=base,inner sep=0pt,text=otMuted,font=\itshape\fontsize{12}{15}\selectfont] at (105mm,129mm) {\matNbActivites{} activités et un guide de l’adulte};
  \pic[scale=3.4] at (105mm,92mm) {opti};
  \node[anchor=base,inner sep=0pt,text=otInk,font=\otMedium\fontsize{10}{12}\selectfont] at (105mm,42mm) {Collection OptiTAB};
  \node[anchor=base,inner sep=0pt,text=otMuted,font=\otLight\fontsize{9}{11}\selectfont] at (105mm,35mm) {\otSite};
\end{scope}
\end{tikzpicture}
\clearpage
% Mentions
\thispagestyle{empty}\vspace*{\fill}
\begin{minipage}{130mm}\fontsize{9}{13}\selectfont\color{otMuted}
{\otSemi\color{otInk}Mathématiques · Maternelle · \matNiveau}\par\medskip
Conception, textes, illustrations et mise en page : collection OptiTAB.\par
\matNbActivites{} activités originales et un guide de l’adulte avec les réponses.\par
Édition du 23 octobre 2026 · Année scolaire 2026-2027.\par\bigskip
Ce cahier suit le programme de l’école maternelle publié au BO n° 41 du
31 octobre 2024, applicable depuis la rentrée 2025, dans sa version consolidée
du BO n° 19 du 7 mai 2026. Il reprend les objectifs que le programme
propose d’aborder \matAge.\par\bigskip
Le programme demande de manipuler avant de passer au papier : chaque fiche
est précédée, dans le guide de l’adulte, d’une activité à faire avec de vrais
objets. Conformément au programme, aucun solide n’est dessiné en
perspective et le nombre zéro n’est pas introduit.\par\bigskip
Toutes les illustrations sont originales. Les quantités, les longueurs et les
réponses de chaque fiche ont été contrôlées par programme informatique, puis
relues page par page.\par\bigskip
© OptiTAB, 2026.\par\otSite
\end{minipage}\clearpage
% Le mot aux adultes
\thispagestyle{matSimple}\phantomsection\pdfbookmark[0]{Le mot aux adultes}{bm-mot}
\matPageTitre{Le mot aux adultes}{Parents, enseignants : comment utiliser ce cahier}
{\fontsize{10.5}{15}\selectfont
À l’école maternelle, les mathématiques se découvrent d’abord en jouant, en
manipulant et en parlant. Ce cahier ne remplace ni le jeu ni la classe : il
prolonge ce que l’enfant a vécu avec de vrais objets et l’aide à fixer ce qu’il
a compris.\par\medskip
\begin{encadrevert}[Cette année, l’enfant va…]\matResume\end{encadrevert}\par\medskip
\begin{encadre}[Une séance réussie, en quatre temps]
\textbf{1. Manipuler.} Faire d’abord l’activité proposée dans le guide de l’adulte avec des jetons, des cubes, des peluches ou les doigts.\par\smallskip
\textbf{2. Lire la consigne.} L’enfant ne sait pas encore lire : l’adulte lit la consigne lentement et montre les pictogrammes.\par\smallskip
\textbf{3. Faire et dire.} L’enfant réalise la fiche et explique ce qu’il fait. Les mots proposés dans le guide l’aident à verbaliser.\par\smallskip
\textbf{4. Valider.} On vérifie ensemble, au besoin en revenant aux objets, puis l’enfant colorie l’étoile de la page.
\end{encadre}\par\medskip
\begin{encadre}[Quelques conseils]
Une ou deux fiches par séance suffisent, pendant une dizaine de minutes. On peut revenir plusieurs fois sur une même notion, à quelques jours d’intervalle. Il n’est pas grave de se tromper : chercher, essayer et recommencer font partie des mathématiques. Valoriser les progrès autant que les réussites aide l’enfant à prendre confiance.
\end{encadre}\par}
\clearpage
% Les pictogrammes
\thispagestyle{matSimple}\phantomsection\pdfbookmark[0]{Les pictogrammes}{bm-pictos}
\matPageTitre{Les pictogrammes}{Pour que l’enfant comprenne ce qu’il doit faire}
{\fontsize{11}{15}\selectfont
\newcommand{\lignePicto}[2]{\par\vspace{3.2mm}\begin{tikzpicture}[baseline=-0.8ex]\fill[white] (0,0) circle (7mm);\draw[otRule,line width=1pt] (0,0) circle (7mm);\pic[scale=.95] at (0,0) {picto #1};\end{tikzpicture}\hspace{6mm}\parbox[c]{130mm}{#2}\par}
\lignePicto{colorier}{\textbf{Colorier} : remplir une forme ou un dessin avec un crayon de couleur.}
\lignePicto{entourer}{\textbf{Entourer} : faire un tour autour de la bonne réponse.}
\lignePicto{relier}{\textbf{Relier} : tracer un trait d’un point à un autre.}
\lignePicto{dessiner}{\textbf{Dessiner} : ajouter des ronds, des objets ou un dessin dans une case.}
\lignePicto{tracer}{\textbf{Tracer} : suivre un chemin ou repasser sur un modèle avec le crayon.}
\lignePicto{barrer}{\textbf{Barrer} : faire une croix sur ce qui ne convient pas.}
\lignePicto{dire}{\textbf{Dire} : expliquer à voix haute ce que l’on voit ou ce que l’on fait.}
\lignePicto{manipuler}{\textbf{Manipuler} : utiliser de vrais objets (cubes, solides, balance…) avant ou pendant la fiche.}
\par\vspace{7mm}
\begin{encadre}[Opti la chouette et l’étoile]
\begin{tikzpicture}[baseline=-1ex]\pic[scale=1.6] at (0,0) {opti};\end{tikzpicture}\hspace{5mm}\parbox[c]{125mm}{Opti la chouette accompagne chaque consigne. En bas de chaque page, une étoile attend d’être coloriée quand l’activité est terminée : l’enfant colorie aussi son étoile sur la page « Mon parcours ».}
\end{encadre}\par}
\clearpage
% Sommaire
\thispagestyle{matSimple}\phantomsection\pdfbookmark[0]{Sommaire}{bm-sommaire}
\matPageTitre{Sommaire}{Maternelle · \matNiveau}
{\fontsize{11}{14}\selectfont\input{sommaire}\par}
\clearpage
% Mon parcours
\thispagestyle{matSimple}\phantomsection\pdfbookmark[0]{Mon parcours}{bm-parcours}
\matPageTitre{Mon parcours}{Quand une page est finie, je colorie son étoile ici aussi}
\vspace{4mm}\input{parcours}
\vfill
\begin{center}\begin{tikzpicture}\pic[scale=2.6] at (0,0) {opti};\node[anchor=west,font=\otMedium\fontsize{13}{17}\selectfont,text=otNavy,text width=110mm] at (2.0,0) {Bravo ! Chaque étoile coloriée est une activité réussie.};\end{tikzpicture}\end{center}
\clearpage
